% ============================================================
% ARCHITECTURE DÉTAILLÉE
% ============================================================
\chapter{Architecture et Conception}

\section{Diagramme de Flux de Données}

Le pipeline de traitement suit un modèle producteur/consommateur avec des threads dédiés pour chaque étape critique. Le flux de données est le suivant :

\begin{enumerate}
    \item La caméra capture des frames à 30 FPS dans un thread dédié.
    \item Les frames sont stockées dans un buffer circulaire thread-safe (\texttt{queue.Queue}).
    \item Le thread principal consomme les frames et les passe au module de pose estimation.
    \item Les landmarks extraits sont transformés en vecteurs de caractéristiques de dimension 78.
    \item La classification est effectuée en parallèle (statique + dynamique) via \texttt{concurrent.futures}.
    \item Les résultats sont affichés localement et diffusés via les interfaces d'accessibilité.
\end{enumerate}

\section{Gestion du Multithreading}

L'architecture exploite le parallélisme à plusieurs niveaux :

\begin{itemize}
    \item \textbf{Thread de capture} : Un thread daemon dédié à l'acquisition vidéo, découplé du traitement pour éviter le blocage.
    \item \textbf{ThreadPoolExecutor} : Utilisation de \texttt{concurrent.futures.ThreadPoolExecutor} avec un nombre de workers égal au nombre de cœurs CPU pour la classification parallèle.
    \item \textbf{Thread TTS} : La synthèse vocale s'exécute dans un thread séparé pour ne pas bloquer le pipeline principal.
    \item \textbf{Thread Flask} : Le serveur MJPEG tourne dans un thread daemon indépendant.
\end{itemize}

\section{Structure du Projet}

\begin{lstlisting}[style=python, caption={Arborescence du projet}]
lsf_translator/
|-- config.py                  # Configuration globale
|-- main.py                    # Pipeline complet (PC)
|-- main_local.py              # Version locale sans reseau
|-- main_pi.py                 # Version Raspberry Pi
|-- train_asl.py               # Entrainement avec dataset ASL
|-- collect_data.py            # Collecte de donnees webcam
|-- requirements.txt           # Dependances PC
|-- requirements_pi.txt        # Dependances Raspberry Pi
|-- models/                    # Modeles TFLite entraines
|   |-- static_mobilenetv3_fp32.tflite
|   |-- static_mobilenetv3_fp16.tflite
|   |-- static_mobilenetv3_int8.tflite
|   |-- labels.npy
|   +-- hand_landmarker.task
+-- modules/
    |-- acquisition.py         # Module 1 : Capture PC
    |-- acquisition_pi.py      # Module 1 : Capture Pi
    |-- pose_estimation.py     # Module 2 : MediaPipe (PC)
    |-- pose_estimation_pi.py  # Module 2 : TFLite (Pi)
    |-- feature_engineering.py # Module 3 : Caracteristiques
    |-- model_training.py      # Module 4 : Entrainement
    +-- accessibility.py       # Module 5 : TTS/Flask/MQTT
\end{lstlisting}

\section{Configuration Globale}

Le fichier \texttt{config.py} centralise tous les paramètres du système :

\begin{lstlisting}[style=python, caption={Paramètres de configuration}]
CAMERA_ID = 0
FRAME_WIDTH = 640
FRAME_HEIGHT = 480
FPS = 30
BUFFER_SIZE = 64

NUM_LANDMARKS = 21
STATIC_FEATURE_DIM = 78
SEQUENCE_LENGTH = 30
MIN_WINDOW = 15
MAX_WINDOW = 45

LSTM_UNITS = 128
BATCH_SIZE = 32
EPOCHS = 50

FLASK_PORT = 5000
MQTT_BROKER = "localhost"
MQTT_PORT = 1883
MQTT_TOPIC = "lsf/predictions"
\end{lstlisting}
